\begin{abstract}
As large language models become the primary interface to computation,
the performance bottleneck shifts from hardware throughput and software
complexity to \emph{meaning itself}: how to represent, compose, and
reason about ideas in a principled way.  We call this the
\textbf{semantic bottleneck}.  This paper introduces \textbf{NEOS}
(Neural Field Operating System), a formal specification that treats the
LLM as a virtual machine and provides an operating-system layer for
cognitive computing.  NEOS is governed by a single master PDE,
$\partial A/\partial t = -\lambda A(x)
  + \alpha\!\int\!K(x,y)\,A(y)\,dy + \iota(x,t)$,
whose three terms---exponential decay, resonance amplification, and
external injection---give rise to attractor formation, coherence
dynamics, and observer-dependent collapse.  The architecture rests on
three pillars: \emph{neural fields} (continuous activation landscapes
as computational substrate), \emph{symbolic reasoning} (a composable
command algebra for field manipulation), and \emph{quantum semantics}
(superposition of meaning with non-commutative, strategy-dependent
collapse).  We validate the framework on two case studies spanning different
domains: (i)~a \emph{software-quality} session in which 69~patterns
were injected over 52~cycles, reaching coherence~0.993, autonomously
expelling one structurally incompatible pattern, and collapsing to five
emergent eigenvectors; and (ii)~a \emph{network-forensics} session in
which 10~patterns injected with default (untuned) parameters over
12~cycles reached coherence~0.93 and autonomously identified an active
AsyncRAT~C2 infection from raw packet-capture data.  The same engine
and master equation produced meaningful structure in both abstract
reasoning and adversarial forensics, providing initial evidence of
domain agility.  To our knowledge, NEOS is the first rigorous
operating-system specification designed for cognitive computing on a
language-model substrate.
\end{abstract}
